
{\bf 04 IV Structure hermitienne}

\subsection{Produit scalaire sur un $\mb{C}$-espace vectoriel} % IV. 2.
\subsubsection{Définition}
\[
h(u,v)=h(v,u)^*\qquad h(u,u)\in\mb{R}^+
\]
\[
h(u,u)=0\ ssi\ u=0\qquad||u||^2=h(u,u)
\]
\subsubsection{Notion d'orthogonalité}
\[
u \perp v \quad alors\quad h(u,v)=0
\]
\[
A^\perp = \{u\in\mc{E},\ \forall a \in A,\ h(a,u)=0\}
\]
famille orthogonale
\[
\{v_i\}\ telle\ que\ \forall i,j,\ i\neq j \Rightarrow h(v_i,v_j)=0
\]

\[
v'_2=\frac{h(v_1,v_2)v_1}{h(v_1,v_1)}\qquad et\qquad v_2-v'_2\ \perp v_1
\]
{\footnotesize\[
h(v_1,v_2-v'_2)=h(v_1,v_2)-h(v_1,v'_2)=
h(v_1,v_2)-h(v_1,\frac{h(v_1,v_2)v_1}{h(v_1,v_1)})=
h(v_1,v_2)-\frac{h(v_1,v_2)}{h(v_1,v_1)}h(v_1,v_1)=0\ \blacksquare
\]}

\subsubsection{Expression d'un vecteur}
$\{e_i\}_{i\in I}$

\[
\forall u\in\mc{E},\ \exists!\{u_i\}_{i\in I}\quad tel\ que\quad u=\sum_iu_ie_i
\]

\[u=
\left(\begin{array}{ c }
 u_1 \\ u_2 \\ \vdots \\ u_n \\
\end{array}\right)
\]

\[
h(u,v)=h(\sum_iu_ie_i,\sum_jv_je_j)=\sum_i\sum_ju_i^*v_jh(e_i,e_j)
\]
$\{e_i\}_{i\in I}$ est une base orthonormé $h(u,v)=\delta_{ij}$
\[
\boxed{h(u,v)=\sum_iu_i^*v_i}
\]

\subsubsection{Notation}
On note 
\[
h(u,v)=\braket{u}{v}
\]

\[
\braket{v}{u}=\braket{u}{v}^*
\]
\subsection{Opérateurs linéaires} % 3.
\subsubsection{Définition} % a.

\newcommand{\wA}[0]{\widehat A}
\[
\begin{array}{ c r c l }
 \wA : & \mc{E} & \to & \mc{E} \\
  & u & \mapsto & v=\wA (u) \\
  &  &  & \ket{v}=\wA \ket{u} \\
\end{array}
\]
\[
\wA \big(\alpha_1\ket{u_1}+\alpha_2\ket{u_2}\big)=
\alpha_1\wA \ket{u_1}+\alpha_2\wA \ket{u_2}
\]
Action de $\wA$
\[
\{\ket{e_i}\}_{i\in I}\qquad\ket{v}=\sum_iv_i\ket{e_i}
\]
\[
\wA\ket{u}=\sum_iu_i\wA\ket{e_i}
\]
\subsubsection{Représentation dans une base orthonormé} % b.
\[
\{\ket{e_i}\}_{i\in I}\qquad h(\ket{e_i},\ket{e_j})=\braket{e_i}{e_j}
\]
\[
\forall i\quad\ket{e_j}\quad\to\quad\wA\ket{e_j}=\sum_ia_{ij}\ket{e_i}\qquad
\]
$\wA$
\[
\wA=\left[\begin{array}{ c c c c }
 a_{11} & a_{12} & \ldots & a_{1n} \\
 a_{21} & a_{22} & \ldots & a_{2n} \\
 \vdots & \vdots &  & \vdots \\
 a_{n1} & a_{n2} & \ldots & a_{nn} \\
\end{array}\right]
\qquad,\qquad\left[\begin{array}{ c }
 a_{11} \\
 a_{21} \\
 \vdots \\
 a_{n1} \\
\end{array}\right]
est\ l'image\ de\ \ket{e_1}\ ...
\]
\subsubsection{Éléments de matrice} % c.
\[
\boxed{a_{ij}=\bra{e_i}\wA\ket{e_j}}
\]

{\footnotesize
\[
\wA\ket{e_j}=\sum_ka_{kj}\ket{e_k}\qquad,\qquad
h(\ket{e_i},\wA\ket{e_j})=a_{ij}=\bra{e_i}\wA\ket{e_j}
\blacksquare\]}
De manière générale, $\bra{u}\wA\ket{v}$ a un sens car le produit matriciel est associatif.
 
\subsubsection{Valeur propre et vecteur propre} % d.
$\ket{\Psi}$ $\lambda$
\[
\wA\ket{\Psi}=\lambda\ket{\Psi}
\]
\[
\lambda\in\mb{C}
\]

\[
\{\ket{}e_i\}_{i\in I}\qquad
\]

\subsection{Opérateur adjoint} % 
\subsubsection{Définition} % a.
Soit 
\[
(\mc{E},h(\cdot,\cdot))
\]
\[
a\in\mc{E}\qquad u\mapsto h(a,u)\quad\in\mb{R} 
\]
Th
\[
\forall f(u) \ lin\acute eaire,\quad\exists!\ b\quad tel\ que\quad
f(u)=\braket{b}{u}
\]

\[
\forall\ket{v}\quad\ket{u}\mapsto h(\ket{v},\wA\ket{u})\quad
est\ lin\acute eaire\ en \ket{u}\ donc
\]
\[
\exists\ket{v'}\quad tel\ que\quad h(\ket{v},\wA\ket{u})=h(\ket{v'},\ket{u})
\]
\[
par\ d\acute efinition,\quad \ket{v'}=\wA^\dagger\ket{v}
\]
Autrement dit,
\[
\forall\ket{v}\quad\exists!\ \ket{v'}\quad tel\ que\quad h(\ket{v},\wA\ket{u})=
h(\ket{v'},\ket{u})
\]
et $\ket{v}$ $\ket{v'}$ $\wA$ $\wA^\dagger$ $\ket{v}\mapsto\ket{v'}$

Autrement dit
\[
\boxed{h(\ket{v},\wA\ket{u})=h(\wA^\dagger\ket{v},\ket{u})}
\]
Autrement dit
\[
\bra{v}\wA\ket{u}=\braket{(\wA^\dagger\ket{v})}{u}
\]

\subsubsection{Représentation matricielle} % b.
$\wA^\dagger$ $\ket{e_i}$
\[
h(\wA^\dagger\ket{e_i},\ket{e_j})=h(\ket{e_i},\wA\ket{e_j})
\]
{\footnotesize
\[
h(\wA^\dagger\ket{e_i},\ket{e_j})=h(\ket{e_j},\wA^\dagger\ket{e_i})^*
=\bra{e_i}\wA^\dagger\ket{e_j}=a^\dagger{ji}
\]
et
\[
h(\ket{e_i},\wA\ket{e_j})=\bra{e_i}\wA\ket{e_j}=a_{ij}
\]
}
d'où
\[
a_{ji}^{\dagger *}=a{ij}\qquad soit\qquad\boxed{a_{ij}^\dagger = a_{ji}^*}
\]

$\wA^\dagger$ $\wA$

